\section{Limitations and Conclusion}\label{sec:conclusion}

\paragraph{Limitations.} Our study is a \emph{single snapshot}, not a time series,
so we observe the state of the storefronts on one day and recover historical
removals only by cross-referencing documented events, not by watching them happen.
It covers only the \emph{Apple} App Store; Google Play, third-party Android stores,
and sideloading are out of scope, so our numbers are a lower bound on app-store
censorship overall. The basket is US-assembled and therefore under-samples
locally important apps, which likely makes some repressive environments look freer
than they are. Availability is an imperfect proxy for intent: as our own data
shows, a missing app can reflect a sanctions-driven developer withdrawal (Candy
Crush in Russia) or a benign regional variant (BBC in the UK) rather than
government censorship, and Iran illustrates a third mechanism entirely, the total
absence of a storefront under sanctions. We mitigate these with control apps,
category-level reasoning, and confidence labels, but we do not claim to resolve
intent perfectly.

\paragraph{Conclusion.} Even within these limits, a single afternoon of public API
queries reproduced every documented removal we tested and surfaced a clean,
interpretable map of app-store censorship: near-total removal in China, a
VPN-focused purge in Russia, a consistent LGBTQ+-app pattern across the Gulf and
Turkey, and a moderate negative correlation with Internet-freedom scores. This
confirms the core premise of the project, that the app-store layer is a real and
measurable censorship chokepoint that network-level tools do not observe, and that
it can be monitored cheaply and safely from outside the censored countries. The
natural next step is to run this pipeline continuously and add Google Play, turning
today's snapshot into the longitudinal observatory the layer still lacks. We
release our code, raw response cache, and availability matrix to make that
straightforward.
